\documentclass{article}
\usepackage[utf8]{inputenc}
\usepackage[T1]{fontenc}
\usepackage{geometry}
\usepackage{graphicx}
\usepackage{float}
\usepackage{hyperref}
\usepackage[english,indonesian]{babel}

\title{Optimasi Pemrosesan Template LaTeX Menggunakan Pendekatan Berbasis AI}
\author{Academic Support Agent}
\date{\today}

\begin{document}
\maketitle

\begin{abstract}

\end{abstract}

\chapter{PENDAHULUAN}


\chapter{TINJAUAN PUSTAKA}


% Original content (commented out):
% \documentclass[a4paper]{article}
\usepackage[utf8]{inputenc}
\usepackage{graphicx}
\usepackage{amsmath}
\usepackage[bottom=2.0cm,top=2.0cm,left=2.0cm,right=2.0cm]{geometry}
\usepackage[portuges]{babel}
\usepackage{indentfirst}
\usepackage{hyperref}  
\hypersetup{colorlinks,citecolor=black,filecolor=black,linkcolor=black,urlcolor=black}
\usepackage[nottoc]{tocbibind}
\usepackage{lipsum}
\usepackage{blindtext}

\begin{document}
\title{Template dos relatórios Física Experimental para Engenharia de Computação}


\begin{titlepage}
	\begin{center}
		\begin{figure}[htb!]
			\begin{flushleft}
				\includegraphics[width=3.9cm]{imagens/LogoCI.jpeg}
			\end{flushleft}
		\end{figure}
        \vspace{-2.5cm}
        \hspace{2.1cm}\Large{\textbf{Universidade Federal da Paraíba}}\\
        \hspace{2.1cm}\Large{Centro de Informática}\\
        \hspace{2.1cm}\large{Física Experimental para Engenharia da Computação}\\
        
        \vspace{200pt}
        
        \LARGE{\textbf{Relatório 01}}\\ 
        \Large{Movimento de Queda Livre via VideoAnálise}\\ 
        
        \vspace{150pt}
        
        \hfill Grupo: XX \\
        
        \vspace{40pt} 
        \hfill Isaac Newton \hspace{20pt} Mat: xxxxxx\\
        \hfill Galileo Galilei \hspace{20pt} Mat: xxxxxx  \\
        \hfill Albert Einstein \hspace{20pt} Mat: xxxxxx\\

        \vspace{25pt}
        \hfill \underline{Professor:}\\
        \hfill Jairo Rocha\\
        
        
        \vspace{\fill}
        \LARGE \bf{\today}
          
	\end{center}
\end{titlepage}

\newpage

%\Large\tableofcontents
%\thispagestyle{empty}
%\newpage
\pagenumbering{arabic}
\large
\lipsum[1]
\section{Resumo}
\lipsum[2-3]
\section{Objetivos}
\lipsum[4-5]
\section{Introdução}
\lipsum[6-8]
\section{Descrição Experimental}
\lipsum[9-10]
\section{Resultados e Análise}
\lipsum[11-12]
\section{Conclusões}
\lipsum[13-15]
\section{Participações}
\lipsum[16]

\bibliographystyle{apalike}
\bibliography{references}




\end{document}


\end{document}